\documentclass[11pt,a4paper,oneside]{report}

\usepackage[utf8]{inputenc}
\usepackage[T1]{fontenc}
\usepackage{lmodern}
\usepackage[margin=1.1in, headheight=16pt, top=1.2in, bottom=1.2in]{geometry}
\usepackage{amsmath,amssymb,amsfonts,amsthm}
\usepackage{booktabs}
\usepackage{graphicx}
\usepackage{tabularx}
\usepackage{hyperref}
\usepackage{fancyhdr}
\usepackage{titlesec}
\usepackage{microtype}
\usepackage{caption}
\usepackage{xcolor}

\definecolor{thesisblue}{RGB}{16, 44, 87}
\definecolor{thesisteal}{RGB}{53, 101, 169}

\hypersetup{
    colorlinks=true,
    linkcolor=thesisblue,
    citecolor=thesisteal,
    urlcolor=thesisteal
}

\pagestyle{fancy}
\fancyhf{}
\fancyhead[L]{\small\slshape\leftmark}
\fancyhead[R]{\small\thepage}
\renewcommand{\headrulewidth}{0.4pt}

% Theorem environments
\newtheorem{theorem}{Theorem}[chapter]
\newtheorem{definition}{Definition}[chapter]
\newtheorem{lemma}[theorem]{Lemma}

% Chapter formatting
\titleformat{\chapter}[display]
  {\normalfont\huge\bfseries\color{thesisblue}}{\chaptertitlename\ \thechapter}{18pt}{\Huge}
\titlespacing*{\chapter}{0pt}{10pt}{30pt}

\begin{document}

% --- TITLE PAGE ---
\begin{titlepage}
    \centering
    \vspace*{1.5cm}
    {\Huge\bfseries\color{thesisblue} Scalable Geometric Deep Learning for Autonomous Multi-Agent Robotics \par}
    \vspace{1.5cm}
    {\Large\itshape Master of Science Thesis in Computer Science \& Engineering \par}
    \vspace{2cm}
    {\large\textbf{Author Name} \par}
    \vspace{0.5cm}
    {\normalsize Department of Electrical Engineering and Computer Science \par}
    {\normalsize University of Excellence \par}
    \vfill
    {\large Supervised by \par}
    {\large\textbf{Prof. Advisor Name} \par}
    \vfill
    {\large\today \par}
\end{titlepage}

% --- ABSTRACT ---
\pagenumbering{roman}
\chapter*{Abstract}
\addcontentsline{toc}{chapter}{Abstract}
Autonomous multi-agent robotic systems deployed in unstructured environments require rigorous distributed state estimation and real-time coordination. In this thesis, we develop a unified geometric deep learning formulation operating over $SE(3)$ manifold representations. We prove convergence bounds for distributed consensus filtering and demonstrate sub-centimeter localization fidelity on embedded edge computing architectures.

\tableofcontents
\listoffigures
\listoftables

\clearpage
\pagenumbering{arabic}

% --- CHAPTER 1 ---
\chapter{Introduction}
\section{Motivation \& Problem Statement}
Multi-robot coordination in GPS-denied environments presents severe mathematical and computational challenges. Classical filtering architectures struggle under non-Gaussian sensory noise and communication latency.

\begin{definition}[$SE(3)$ Lie Group Pose]
A rigid body pose $T \in SE(3)$ consists of a rotation matrix $R \in SO(3)$ and a translation vector $p \in \mathbb{R}^3$, represented in homogeneous coordinates as:
\begin{equation}
T = \begin{bmatrix} R & p \\ 0 & 1 \end{bmatrix} \in \mathbb{R}^{4 \times 4}
\end{equation}
\end{definition}

\chapter{Theoretical Foundations}
\section{Manifold Consensus Optimization}
Let $\mathcal{G} = (\mathcal{V}, \mathcal{E})$ denote the inter-robot communication topology graph. The distributed optimization objective is formulated as:
\begin{equation}
\min_{x_i \in \mathcal{M}} \sum_{i \in \mathcal{V}} f_i(x_i) + \sum_{(i, j) \in \mathcal{E}} d_{\mathcal{M}}^2(x_i, x_j)
\end{equation}
where $d_{\mathcal{M}}(\cdot, \cdot)$ denotes the Riemannian geodesic distance metric on manifold $\mathcal{M}$.

\begin{theorem}[Geometric Convergence Bound]
Under strongly geodesic convexity on Riemannian manifolds with sectional curvature bounded by $\kappa$, the consensus algorithm achieves linear convergence $\mathcal{O}(\gamma^k)$ with rate $\gamma < 1$.
\end{theorem}

\chapter{Experimental Results \& Conclusion}
\begin{table}[h!]
\centering
\caption{Benchmark Trajectory Accuracy Across Heterogeneous Environments}
\begin{tabular}{lccc}
\toprule
\textbf{Algorithm Architecture} & \textbf{ATE (m)} & \textbf{RPE (deg/m)} & \textbf{Latency (ms)} \\
\midrule
Standard EKF SLAM               & 0.284 & 1.42 & \textbf{1.8} \\
Factor Graph Smoothing          & 0.086 & 0.45 & 12.4 \\
\textbf{Ours (Manifold GNN)}    & \textbf{0.019} & \textbf{0.09} & 4.2 \\
\bottomrule
\end{tabular}
\end{table}

\section{Summary \& Future Work}
This work established a scalable framework bridging Riemannian geometry and graph neural networks for robotic swarms. Future directions include hardware validation under non-line-of-sight RF environments.

\end{document}
